\documentclass[prx,amsmath,amssymb,superscriptaddress,longbibliography]{revtex4-2}
% \documentclass[10pt,prl,superscriptaddress,longbibliography]{revtex4-2}
\usepackage{xcolor}
\definecolor{BLACK}{rgb}{0,0,0}
\usepackage{bm}
\usepackage{amsmath}
\usepackage[colorlinks=true,linkcolor=blue,anchorcolor=blue,citecolor=blue,urlcolor=blue]{hyperref}


\usepackage{booktabs}
\usepackage{mathrsfs}
\makeatletter
\def\ltxu@dotsep{0.5em}
\makeatother
\newcommand{\Tr}{\operatorname{Tr}}
\newcommand{\Rea}{\operatorname{Re}}
\newcommand{\Ima}{\operatorname{Im}}
\newcommand{\dd}{\mathrm{d}}
\newcommand{\ii}{\mathrm{i}}
\newcommand{\symab}{\mathop{\mathrm{Sym}}\nolimits_{ab}}
\newcommand{\code}[1]{\texttt{\detokenize{#1}}}
\allowdisplaybreaks[2]
\begin{document}
\title{Personally implemented formulae in WannierBerri\\
\color{blue}Working handbook: geometry, covariant derivatives, and interpolation}
\author{Jiayu Li}
\email{jyli93@hku.hk}
\affiliation{Department of Physics, Southern University of Science and Technology, Shenzhen, China}
\affiliation{New Cornerstone Science Laboratory, Department of Physics, The University of Hong Kong, Hong Kong, China}
\date{\textcolor{blue}{Working copy, 4 October 2026}}
\maketitle
\color{blue}
\noindent\textbf{Revision colours.} Black text is retained from the original
\code{formulae.tex}; blue text marks additions and rewritten passages.
Unchanged equations remain black even when labels have been added.


This working copy consolidates the results in \code{formulae.tex},
\code{Wannier.tex}, and \code{QGD.tex}.
It establishes common conventions for the subsequent reconciliation of the
v1.7 and v1.8 calculators. Formula identities, representation assumptions,
and implementation status are stated separately.
The original source notes remain unchanged. Source corrections and unfinished
results are recorded in Appendix~\ref{app:corrections}.

% A compact contents list keeps the overview on the title page.
\begingroup
\small
\makeatletter
\def\l@f@section{\addvspace{0.4em}}
\makeatother
\tableofcontents
\endgroup
\newpage

\section{\textcolor{black}{Basic ingredients}}
\label{sec:basic}

\subsection{Notation, units, and scope}
\label{sec:notation}
We use Cartesian wave-vector derivatives
$\partial_a=\partial/\partial k_a$ and $\epsilon_{xyz}=+1$.
Repeated Cartesian indices are summed only in explicitly contracted expressions;
free indices remain free. An isolated band is denoted by $n$.
In block formulae, $m,n\in A$ and $l,p\in B$, where the two groups
are separated by a nonzero gap at the wave vector under consideration.
Degeneracies within a group are allowed. Primed band indices are summed over
their indicated group.
\begin{equation}
 \Delta_{rs}=\varepsilon_r-\varepsilon_s,\qquad
 \int[\dd\bm{k}]=\int_{\mathrm{BZ}}\frac{\dd^d k}{(2\pi)^d},\qquad
 e>0,\quad q_{\mathrm{electron}}=-e.
 \label{eq:conventions}
\end{equation}
The geometric expressions use energies and wave vectors, without silently
setting $\hbar=1$. In particular,
\begin{equation}
 v_n^a=\frac{1}{\hbar}\partial_a\varepsilon_n,\qquad
 E_n^{a}=\partial_a\varepsilon_n.
 \label{eq:velocity-units}
\end{equation}
If energy is in eV and length in \AA, $E_n^a$ is in eV\,\AA;
it is not yet a velocity in m/s. The conversion is
$v_n^a=(1\,\mathrm{eV\ in\ J})(1\,\text{\AA\ in m})E_n^a/\hbar$.

\begin{center}
\setlength{\tabcolsep}{6pt}
\begin{tabular}{lll}
\toprule
Symbol & Meaning & Relation to the source notes\\
\midrule
$Q_n^{ab}$ & quantum geometric tensor (QGT) & QG; called QMT in \code{formulae.tex}\\
$G_n^{ab}$ & quantum metric & $g_n^{ab}$ in the other two notes\\
$\mathcal G_n^{ab}$ & Berry connection polarizability (BCP) & original normalization retained\\
$D^a$ & $U^\dagger\partial_a U$ & coefficient-space connection\\
$\bar A^a$ & $U^\dagger A^{W,a}U$ & rotated Wannier connection\\
$M^a$ & $\bar A^a+\ii D^a$ & auxiliary connection in BCP formulae\\
$\widetilde G^{ab}$ & energy-weighted geometric matrix & distinct from the metric $G_n^{ab}$\\
$\mathfrak m^a$ & interband magnetic kernel & separates magnetic $M$ from BCP $M$\\
\bottomrule
\end{tabular}
\end{center}

{\color{black}
\subsection{Berry connection}

Single band Berry connection (Abelian Berry connection):
\begin{equation}
    A_n^i(\bm{k}) = i \langle u_n(\bm{k})| \partial_{k_i} u_n(\bm{k}) \rangle \equiv i \langle{u_n}| \partial_{i}{u_n}\rangle 
\label{eq:berry-connection}\end{equation}
Interband Berry connection (non-Abelian Berry connection):
\begin{equation}
    A_{nm}^i(\bm{k}) = i \langle u_n(\bm{k})| \partial_{k_i} u_m(\bm{k}) \rangle \equiv i \langle{u_n}| \partial_{i}{u_m}\rangle 
\end{equation}

}
The matrix Berry connection is Hermitian, $(A^a)^\dagger=A^a$.
Its full matrix transforms inhomogeneously under a wave-vector-dependent
change of basis. Off-block elements transform covariantly under independent
rotations within $A$ and $B$.

\subsection{Quantum geometric tensor (QGT)}
{\color{black}
Quantum geometry tensor: 
\begin{align}
     Q_{n}^{ij}(\bm{k}) &=  \langle{\partial_{k_i} u_n(\bm{k})}|(1 - |u_n(\bm{k})\rangle\langle{u_n(\bm{k})}|)|{\partial_{k_j} u_n(\bm{k})}\rangle \nonumber \\
     &\equiv \langle \partial_{i} u_n |(1 - |u_n\rangle\langle u_n |)|{\partial_{j} u_n}\rangle  \nonumber  \\
     &= \sum_{m(\neq n)}  \langle \partial_i u_n | u_m \rangle\langle u_m | \partial_j u_n \rangle \nonumber \\
     &= \sum_{m(\neq n)}  A_{nm}^i A_{mn}^j \nonumber \\
     &= G_n^{ij}(\bm{k}) - \frac{i}{2}\Omega_n^{ij}(\bm{k}) 
\label{eq:qgt}\end{align}
Quantum metric tensor: 
\begin{equation}
    G_{n}^{ij}(\bm{k}) = \mathrm{Re}~Q_{n}^{ij}(\bm{k}) = \mathrm{Re} \sum_{m(\neq n)}  A_{nm}^i A_{mn}^j
\label{eq:metric-curvature}\end{equation}
Berry curvature tensor:
\begin{equation}
    \Omega_{n}^{ij}(\bm{k}) = -2\mathrm{Im}~Q_{n}^{ij}(\bm{k}) = -2\mathrm{Im} \sum_{m(\neq n)}  A_{nm}^i A_{mn}^j
\end{equation}
Berry curvature vector:
\begin{equation}
    \Omega_{n}^{i}(\bm{k}) = \frac{1}{2}\epsilon^{ijk}\Omega_{n}^{jk}(\bm{k}) = -\epsilon^{ijk}\mathrm{Im} \sum_{m(\neq n)}  A_{nm}^j A_{mn}^k = \epsilon^{ijk} \partial_j A_n^k
\label{eq:curvature-vector}\end{equation}
}
The band sum in Eq.~\eqref{eq:qgt} requires a complete complement to $n$.
A finite-model band sum gives the geometry of that model; missing states in a
Wannier representation require the external terms of Sec.~\ref{sec:interpolation}.
The metric is symmetric and positive semidefinite, whereas
$\Omega_n^{ab}=-\Omega_n^{ba}$.

{\color{black}
\subsection{Berry connection polarizability (BCP)}

Berry connection polarizability~\cite{GaoPRL2014FieldInducedPositional}: 
\begin{equation}
    \mathcal{G}_n^{ij}(\bm{k}) = \frac{1}{2} \sum_{m(\neq n)}  \frac{A_{nm}^i A_{mn}^j + A_{nm}^j A_{mn}^i}{\varepsilon_n - \varepsilon_m} = \mathrm{Re}\sum_{m(\neq n)}  \frac{A_{nm}^i A_{mn}^j}{\varepsilon_n - \varepsilon_m}
\label{eq:bcp-single}\end{equation}

}
The normalization in Eq.~\eqref{eq:bcp-single} is retained from the original notes.
Any response formula imported from another convention must first be converted
to this definition; a factor of two must not be inferred from the name BCP alone.
The units of $Q$, $G$, and $\Omega$ are length squared; those of
$\mathcal G$ are length squared divided by energy.

\section{Covariant derivatives and Wannier conventions}
\label{sec:covariant}

\subsection{Wannier and Hamiltonian gauges}
With Wannier centers $\bm t_i$, define
\begin{align}
 |u_{j\bm k}^{W}\rangle
 &=\sum_{\bm R}e^{\ii\bm k\cdot(\bm R+\bm t_j-\bm r)}|\bm Rj\rangle,
 \\
 \mathcal H_{ij}^{W}(\bm k)
 &=\sum_{\bm R}e^{\ii\bm k\cdot\bm R_{ij}}
 \langle\bm0 i|\hat{\mathcal H}|\bm Rj\rangle,
 \qquad \bm R_{ij}=\bm R+\bm t_j-\bm t_i,
 \label{eq:fourier}\\
 \mathcal H^W U&=U\varepsilon,\qquad
 |u_n\rangle=\sum_j|u_j^W\rangle U_{jn},\qquad
 \bar X=U^\dagger X^W U.
 \label{eq:gauges}
\end{align}
Normalization of the Bloch sums is understood.
Fourier derivatives follow by multiplying every real-space term by
$\ii R_{ij}^a$, with one factor for each derivative.

For the coefficient vectors (columns of $U$),
\begin{align}
 D^a&=U^\dagger\partial_a U,\qquad
 (D^a)^\dagger=-D^a,\qquad (D_{nm}^a)^*=-D_{mn}^a,
 \label{eq:D}\\
 D_{ln}^a&=\frac{\mathcal H_{ln}^{,a}}{\varepsilon_n-\varepsilon_l},
 \quad l\in B,\ n\in A,\qquad
 \mathcal H^{,a}=U^\dagger(\partial_a\mathcal H^W)U,
 \label{eq:D-gap}\\
 |\partial_a u_n\rangle
 &=\sum_j|\partial_a u_j^W\rangle U_{jn}+\sum_r|u_r\rangle D_{rn}^a,
 \\
 A^{H,a}&=\bar A^a+\ii D^a\equiv M^a,\qquad
 A_{ij}^{W,a}=\ii\langle u_i^W|\partial_a u_j^W\rangle.
 \label{eq:M}
\end{align}
Diagonal entries of $D$ depend on the phase choice.
Equation~\eqref{eq:D-gap} is used only across separated blocks, never to divide
by a zero gap within a degenerate group.
In general, $D_{nm}^a\ne\langle u_n|\partial_a u_m\rangle$:
the latter also contains the derivative of the Wannier basis.

\subsection{Comma, ordinary, and covariant derivatives}
\label{sec:derivative-conventions}
For a specified Wannier-gauge matrix $X^W$, the comma derivative means
\begin{equation}
 X^{,d}=U^\dagger(\partial_d X^W)U,\qquad
 \partial_d\bar X=X^{,d}+\bar X D^d-D^d\bar X.
 \label{eq:ordinary-derivative}
\end{equation}
Thus $X^{,d}$ is not the ordinary derivative of $\bar X$.
In a fixed basis for an abstract operator it becomes
$\langle n|\partial_d\hat X|m\rangle$.
Such an operator-derivative matrix is covariant; it is the ordinary derivative
of a transformed matrix that acquires derivatives of the gauge rotation.

Under $U\mapsto U\,\mathrm{diag}(W_A,W_B)$,
$X_{AB}\mapsto W_A^\dagger X_{AB}W_B$.
Using the coefficient-space connection, define~\cite{Liu2023CovariantderivativesBerry}
\begin{equation}
 X_{CD}^{:d}=\partial_d X_{CD}+D_{CC}^d X_{CD}-X_{CD}D_{DD}^d,
 \quad C,D\in\{A,B\}.
 \label{eq:cov-derivative}
\end{equation}
Here $C,D$ label blocks, while the superscript $d$ labels a Cartesian direction.
The symbol $X$ in this equation denotes a Hamiltonian-gauge matrix.
For a rotated Wannier matrix the resulting block rules are
\begin{align}
 X_{AB}^{:d}&=X_{AB}^{,d}-D_{AB}^dX_{BB}+X_{AA}D_{AB}^d,\label{eq:block-ab}\\
 X_{BA}^{:d}&=X_{BA}^{,d}-D_{BA}^dX_{AA}+X_{BB}D_{BA}^d,\\
 X_{AA}^{:d}&=X_{AA}^{,d}-D_{AB}^dX_{BA}+X_{AB}D_{BA}^d.
 \label{eq:block-aa}
\end{align}
For the Hamiltonian, $\mathcal H_{AA}^{:d}=\mathcal H_{AA}^{,d}$ and
$\mathcal H_{AB}^{:d}=0$.
These statements concern rotations within a fixed Wannier representation.
A connection defined using the full physical Bloch basis must be used
consistently if that representation is changed.

\subsection{Trace, product, and chain rules}
The trace and product rules eliminate gauge-dependent intrablock connections:
\begin{align}
 \partial_d\Tr_A X&=\Tr_A X^{:d},
 \label{eq:trace}\\
 (X_{CD}Y_{DE})^{:d}&=X_{CD}^{:d}Y_{DE}+X_{CD}Y_{DE}^{:d}.
 \label{eq:product}
\end{align}
For $\widetilde f=f(\mathcal H)$, evaluation in the energy eigenbasis gives
\begin{equation}
 \widetilde f_{mn}^{:d}=\mathcal H_{mn}^{:d}
 \begin{cases}
 f'(\varepsilon_n),&\varepsilon_m=\varepsilon_n,\\[2pt]
 \displaystyle\frac{f(\varepsilon_m)-f(\varepsilon_n)}{\varepsilon_m-\varepsilon_n},
 &\varepsilon_m\ne\varepsilon_n.
 \end{cases}
 \label{eq:chain}
\end{equation}
The divided difference is essential when a block contains distinct energies.
At zero temperature, apply the smooth-function rule before taking the
appropriate distributional limit of the occupation.

For compact second derivatives, define the bilinear product expansion
\begin{align}
 \mathscr P_d(X,Y)&=X^{:d}Y+XY^{:d},\nonumber\\
 \mathscr P_{de}(X,Y)&=X^{:de}Y+X^{:d}Y^{:e}
                 +X^{:e}Y^{:d}+XY^{:de},
 \label{eq:second-product}\\
 X^{:de}&\equiv(X^{:d})^{:e}.
 \nonumber
\end{align}
Matrix order is preserved. Ordered covariant derivatives need not commute;
ordinary derivatives of a smooth traced scalar do commute.

\section{Gauge-covariant geometry and Wannier interpolation}
\label{sec:interpolation}

\subsection{Projectors and geometric matrices}
Let $P=\sum_{n\in A}|u_n\rangle\langle u_n|$ and $Q=1-P$.
For physical Bloch states, $Q$ includes states outside a finite Wannier space.
The projected geometric matrices are
\begin{align}
 \widetilde F_{mn}^{ab}
 &=\langle\partial_a u_m|Q|\partial_b u_n\rangle,
 \label{eq:Fprojected}\\
 \widetilde G^{ab}&=\tfrac12(\mathcal H_{AA}\widetilde F^{ab}
                            +\widetilde F^{ab}\mathcal H_{AA}),
 \label{eq:Gweighted}\\
 \widetilde H_{mn}^{ab}
 &=\langle\partial_a u_m|Q\hat{\mathcal H}Q|\partial_b u_n\rangle.
 \label{eq:Hprojected}
\end{align}
Each transforms as $\widetilde O\mapsto W_A^\dagger\widetilde O W_A$ and obeys
$\widetilde O_{mn}^{ab}=(\widetilde O_{nm}^{ba})^*$.
In the Hamiltonian gauge,
$\widetilde G_{mn}^{ab}=(\varepsilon_m+\varepsilon_n)\widetilde F_{mn}^{ab}/2$.
The metric and curvature matrices are
\begin{equation}
 \widetilde g^{ab}=\tfrac12(\widetilde F^{ab}+\widetilde F^{ba}),\qquad
 \widetilde\Omega^c=\ii\epsilon_{cab}\widetilde F^{ab}.
 \label{eq:projected-geometry}
\end{equation}
For $A=\{n\}$ these recover Sec.~\ref{sec:basic}.
For a larger $A$, $\Tr_A\widetilde g$ is the metric of the whole subspace;
it is generally not $\sum_{n\in A}G_n$.

For a continuum Hamiltonian with kinetic energy $(\bm p+\hbar\bm k)^2/(2m_e)$
and a local, momentum-independent potential, the fourth matrix in
\code{Wannier.tex} is
\begin{equation}
 \widetilde S_{mn}^{ab}=\frac{\hbar^2}{m_e}\delta_{ab}\delta_{mn}
                  +2(\widetilde G_{mn}^{ab}-\widetilde H_{mn}^{ab}).
 \label{eq:S-continuum}
\end{equation}
For $A=\{n\}$, $\Rea\widetilde S_{nn}^{ab}=\partial_a\partial_b\varepsilon_n$.
The free-electron term cannot be substituted for the bare Hessian of a
general tight-binding or nonlocal Hamiltonian; use Sec.~\ref{sec:energy} there.

\subsection{Internal and external terms}
Define Wannier-gauge matrices
\begin{align}
 F_{ij}^{W,ab}&=\langle\partial_a u_i^W|\partial_b u_j^W\rangle,\nonumber\\
 B_{ij}^{W,a}&=\ii\langle u_i^W|\hat{\mathcal H}|\partial_a u_j^W\rangle,\nonumber\\
 H_{ij}^{W,ab}&=\langle\partial_a u_i^W|\hat{\mathcal H}|\partial_b u_j^W\rangle.
 \label{eq:external-data}
\end{align}
The two-index $H^{W,ab}$ is an energy-weighted geometric matrix,
not $\partial_a\partial_b\mathcal H^W$.
Then
\begin{align}
 \widetilde F^{ab}_{\mathrm{int}}
 &=-D_{AB}^aD_{BA}^b,\label{eq:Fint}\\
 \widetilde F^{ab}_{\mathrm{ext}}
 &=\bar F_{AA}^{ab}-\bar A_{AA}^a\bar A_{AA}^b
    +\ii D_{AB}^a\bar A_{BA}^b+\ii\bar A_{AB}^a D_{BA}^b,\label{eq:Fext}\\
 \widetilde F^{ab}&=\widetilde F^{ab}_{\mathrm{int}}+
                    \widetilde F^{ab}_{\mathrm{ext}}.
\end{align}
Here $B$ is the complement of $A$ within the represented Wannier space;
$\bar F$ retains the additional basis-derivative information.
Keeping only $\widetilde F_{\mathrm{int}}$ corresponds to the fixed-basis
Hamiltonian geometry.

For $\widetilde H$, assuming the represented active states are eigenstates
of the physical Hamiltonian (or lie in the exact frozen subspace),
\begin{align}
 \widetilde H^{ab}_{\mathrm{int}}&=-D_{AB}^a\mathcal H_{BB}D_{BA}^b,\nonumber\\
 \widetilde H^{ab}_{\mathrm{ext}}
 &=\bar H_{AA}^{ab}-\bar A_{AA}^a\mathcal H_{AA}\bar A_{AA}^b
    +\ii D_{AB}^a\bar B_{BA}^b
    +\ii(\bar B^a)^\dagger_{AB}D_{BA}^b.
 \label{eq:Hinterpolation}
\end{align}
Outside that exact subspace, projection and disentanglement errors must be
assessed before interpreting the interpolation as a physical matrix element.

A useful sign check is
\begin{equation}
 \Omega_{A,\mathrm{int}}^{ab}
 =2\Ima\Tr_A(D_{AB}^aD_{BA}^b),\qquad
 \widetilde\Omega^z=\ii(\widetilde F^{xy}-\widetilde F^{yx}).
 \label{eq:curvature-check}
\end{equation}

\subsection{First and second geometric derivatives}
The internal first derivative is
\begin{equation}
 \widetilde F_{\mathrm{int}}^{ab:d}
 =-D_{AB}^{a:d}D_{BA}^b-D_{AB}^aD_{BA}^{b:d}.
 \label{eq:Ffirst}
\end{equation}
The external first derivative contains all seven product-rule terms:
\begin{align}
 \widetilde F_{\mathrm{ext}}^{ab:d}
 &=\bar F_{AA}^{ab:d}
 -\bar A_{AA}^{a:d}\bar A_{AA}^b-\bar A_{AA}^a\bar A_{AA}^{b:d}
 \nonumber\\
 &\quad+\ii D_{AB}^{a:d}\bar A_{BA}^b+\ii D_{AB}^a\bar A_{BA}^{b:d}
 +\ii\bar A_{AB}^{a:d}D_{BA}^b+\ii\bar A_{AB}^aD_{BA}^{b:d}.
 \label{eq:Fext-first}
\end{align}
All second-order terms are retained compactly by Eq.~\eqref{eq:second-product}:
\begin{align}
 \widetilde F_{\mathrm{int}}^{ab:de}
 &=-\mathscr P_{de}(D_{AB}^a,D_{BA}^b),\nonumber\\
 \widetilde F_{\mathrm{ext}}^{ab:de}
 &=\bar F_{AA}^{ab:de}-\mathscr P_{de}(\bar A_{AA}^a,\bar A_{AA}^b)
 \nonumber\\
 &\quad+\ii\mathscr P_{de}(D_{AB}^a,\bar A_{BA}^b)
       +\ii\mathscr P_{de}(\bar A_{AB}^a,D_{BA}^b).
 \label{eq:Fsecond}
\end{align}
Consequently, for $s=d$ or $de$,
\begin{equation}
 \widetilde g^{ab:s}=\tfrac12(\widetilde F^{ab:s}+\widetilde F^{ba:s}),\qquad
 \widetilde\Omega^{c:s}=\ii\epsilon_{cab}\widetilde F^{ab:s}.
 \label{eq:geometry-derivatives}
\end{equation}
Taking $\Tr_A$ gives the ordinary first or second derivative of the
corresponding subspace trace.
The same product rule differentiates Eqs.~\eqref{eq:Gweighted} and
\eqref{eq:Hinterpolation}; derivatives of every Hamiltonian factor are required.

\section{BCP and energy-denominator derivatives}
\label{sec:bcp}

\subsection{A symmetric block construction}
For off-block elements define $M$ by Eq.~\eqref{eq:M} and
\begin{equation}
 N_{rs}^a=\frac{M_{rs}^a}{\Delta_{rs}},\qquad
 (M_{rs}^a)^*=M_{sr}^a,\qquad (N_{rs}^a)^*=-N_{sr}^a.
 \label{eq:N}
\end{equation}
The finite-model energy-weighted tensor and BCP are
\begin{align}
 \widetilde{\mathcal F}^{ab}
 &=\tfrac12(N_{AB}^a M_{BA}^b-M_{AB}^a N_{BA}^b),\label{eq:weightedF}\\
 \widetilde{\mathcal G}^{ab}
 &=\tfrac12(\widetilde{\mathcal F}^{ab}+\widetilde{\mathcal F}^{ba}).
 \label{eq:bcp-block}
\end{align}
In the energy eigenbasis,
\begin{equation}
 \widetilde{\mathcal F}_{mn}^{ab}
 =\frac12\sum_{l\in B}M_{ml}^a M_{ln}^b
 \left(\frac1{\Delta_{ml}}+\frac1{\Delta_{nl}}\right).
 \label{eq:weightedF-components}
\end{equation}
For a single active band and a complete model complement,
$\widetilde{\mathcal G}_{nn}^{ab}=\mathcal G_n^{ab}$.
A calculation using only $M$ within a truncated Wannier space does not restore
energy-denominator contributions from all excluded physical bands.

Let $T^{ab}=N_{AB}^aM_{BA}^b-M_{AB}^aN_{BA}^b$.
Hermiticity gives $(T^{ab})^\dagger=T^{ba}$, so the code-friendly form is
\begin{equation}
 \widetilde{\mathcal G}^{ab}=\tfrac14[T^{ab}+(T^{ab})^\dagger].
 \label{eq:bcp-code}
\end{equation}
This explains the division by four after band-index Hermitian symmetrization
in the archived \code{BCP} class.

\subsection{Connection derivatives and the general block rule}
The first derivative of the off-block coefficient connection is
\begin{align}
 D_{ln}^{a:d}=\frac{1}{\varepsilon_n-\varepsilon_l}
 \big[\mathcal H_{ln}^{,ad}
 &+\mathcal H_{lp}^{,a}D_{pn}^d-D_{lm}^d\mathcal H_{mn}^{,a}
 \nonumber\\[-2pt]
 &+\mathcal H_{lp}^{,d}D_{pn}^a-D_{lm}^a\mathcal H_{mn}^{,d}\big],
 \label{eq:Dfirst}
\end{align}
where $m\in A$ and $p\in B$ are summed.
Also $M^{a:d}=\bar A^{a:d}+\ii D^{a:d}$, with the corresponding relation at second order.

To differentiate an energy denominator covariantly, define the Sylvester operator
\begin{equation}
 \mathscr L_{CD}(Y)=\mathcal H_{CC}Y-Y\mathcal H_{DD}.
 \label{eq:sylvester}
\end{equation}
Its inverse exists for separated spectra. Equation~\eqref{eq:N} is the
Hamiltonian-gauge evaluation of $\mathscr L_{CD}(N^a)=M^a$.
For $\mathscr L_{CD}(Y)=R$, product differentiation gives
\begin{align}
 \mathscr L_{CD}(Y^{:d})
 &=R^{:d}-\mathcal H_{CC}^{:d}Y+Y\mathcal H_{DD}^{:d},\label{eq:sylvester1}\\
 \mathscr L_{CD}(Y^{:de})
 &=R^{:de}-\mathcal H_{CC}^{:de}Y
 -\mathcal H_{CC}^{:d}Y^{:e}-\mathcal H_{CC}^{:e}Y^{:d}
 \nonumber\\
 &\quad+Y^{:d}\mathcal H_{DD}^{:e}+Y^{:e}\mathcal H_{DD}^{:d}
       +Y\mathcal H_{DD}^{:de}.
 \label{eq:sylvester2}
\end{align}
Set $(Y,R)=(N^a,M^a)$ for BCP denominators, or
$(Y,R)=(D^a_{BA},-\mathcal H_{BA}^{,a})$ for coefficient connections.
These identities preserve intrablock matrix products and supply a general
reference for auditing the multiband implementation.

\subsection{Isolated-band quotient formulae}
For individual nondegenerate eigenlines, write the phase-covariant derivative
as $;d$, to distinguish it from a multiband block derivative.
Let $\Delta_{rs,d}=E_r^d-E_s^d$ and
$\Delta_{rs,de}=\partial_e\partial_d(\varepsilon_r-\varepsilon_s)$.
Then
\begin{align}
 N_{rs}^{a;d}
 &=\frac{M_{rs}^{a;d}}{\Delta_{rs}}
   -\frac{M_{rs}^{a}\Delta_{rs,d}}{\Delta_{rs}^2},\label{eq:Nfirst}\\
 N_{rs}^{a;de}
 &=\frac{M_{rs}^{a;de}}{\Delta_{rs}}
   -\frac{M_{rs}^{a;d}\Delta_{rs,e}+M_{rs}^{a;e}\Delta_{rs,d}
           +M_{rs}^a\Delta_{rs,de}}{\Delta_{rs}^2}
   +\frac{2M_{rs}^a\Delta_{rs,d}\Delta_{rs,e}}{\Delta_{rs}^3}.
 \label{eq:Nsecond}
\end{align}
These are the compact quotient formulae in \code{Wannier.tex}.
They must not replace Eqs.~\eqref{eq:sylvester1}--\eqref{eq:sylvester2}
for arbitrary multiband rotations. The archived \code{DerBCP} uses diagonal
velocity differences; its general-block equivalence remains a reconciliation item.

\subsection{BCP dipole and quadrupole kernels}
With either consistently evaluated block derivatives or isolated-band derivatives,
\begin{align}
 \widetilde{\mathcal F}^{ab:d}
 &=\tfrac12\big[
 N_{AB}^{a:d}M_{BA}^b+N_{AB}^aM_{BA}^{b:d}
 -M_{AB}^{a:d}N_{BA}^b-M_{AB}^aN_{BA}^{b:d}\big],\label{eq:bcp-first}\\
 \widetilde{\mathcal F}^{ab:de}
 &=\tfrac12\big[\mathscr P_{de}(N_{AB}^a,M_{BA}^b)
                  -\mathscr P_{de}(M_{AB}^a,N_{BA}^b)\big],\label{eq:bcp-second}\\
 \widetilde{\mathcal G}^{ab:s}
 &=\tfrac12(\widetilde{\mathcal F}^{ab:s}
             +\widetilde{\mathcal F}^{ba:s}),\qquad s=d,de.
 \label{eq:bcp-sym}
\end{align}
The two factors of $1/2$ in this construction are independent.
For a smooth subspace, $\Tr_A\widetilde{\mathcal G}^{ab:d}$ and
$\Tr_A\widetilde{\mathcal G}^{ab:de}$ are the dipole and quadrupole
geometric kernels before occupations and response prefactors are applied.

\section{Explicit band-sum and matrix formulae from QGD}
\label{sec:qgd}

\subsection{Domain and matrix ingredients}
This section applies to a smooth Hamiltonian in a fixed orthonormal basis,
with nondegenerate bands and nonzero denominators.
It gives an independent representation of the internal geometry.
Define energy-derivative matrices, rather than physical velocities,
\begin{align}
 V_{nm}^a&=\langle n|\partial_a\mathcal H|m\rangle,
 &\Gamma_{nm}^{ab}&=\langle n|\partial_a\partial_b\mathcal H|m\rangle,
 \nonumber\\
 (K_q^a)_{nm}&=\begin{cases}V_{nm}^a/\Delta_{nm}^q,&n\ne m,\\0,&n=m,\end{cases}
 &(V_F^a)_{nm}&=\delta_{nm}E_n^a.
 \label{eq:qgd-ingredients}
\end{align}
Division by $\Delta^q$ is elementwise and excludes the diagonal.
Products of $K_q$ are ordinary matrix products.
In this convention $D_{nm}^a=-(K_1^a)_{nm}$ and
$A_{nm}^a=-\ii(K_1^a)_{nm}$ for $n\ne m$.
Thus $K_1$ alone is not the interband Berry connection.

The QGT can be evaluated by either expression
\begin{equation}
 Q_n^{ab}=\sum_{m\ne n}\frac{V_{nm}^aV_{mn}^b}{\Delta_{nm}^2}
          =-[K_1^aK_1^b]_{nn}.
 \label{eq:qgt-matrix}
\end{equation}
Only the diagonal of the matrix product is asserted to equal the
single-band QGT. Its off-diagonal entries are not the general projected
matrix of Eq.~\eqref{eq:Fprojected}.

\subsection{QGD-I and QGD-II}
Write $\partial_cQ_n^{ab}=\mathcal I_n^{ab;c}+\mathcal J_n^{ab;c}$.
The one-sum contribution is
\begin{equation}
 \mathcal I_n^{ab;c}=\sum_{m\ne n}\left[
 \frac{\Gamma_{nm}^{ca}V_{mn}^b+V_{nm}^a\Gamma_{mn}^{cb}}{\Delta_{nm}^2}
 -\frac{2V_{nm}^aV_{mn}^b(E_n^c-E_m^c)}{\Delta_{nm}^3}\right].
 \label{eq:qgd-I}
\end{equation}
The two-sum contribution is
\begin{align}
 \mathcal J_n^{ab;c}=\sum_{m\ne n}\bigg[
 &\sum_{p\ne n}\frac{V_{np}^cV_{pm}^aV_{mn}^b
                        +V_{nm}^aV_{mp}^bV_{pn}^c}
                       {\Delta_{nm}^2\Delta_{np}}
 \nonumber\\
 +&\sum_{p\ne m}\frac{V_{np}^aV_{pm}^cV_{mn}^b
                        +V_{nm}^aV_{mp}^cV_{pn}^b}
                       {\Delta_{nm}^2\Delta_{mp}}\bigg].
 \label{eq:qgd-II}
\end{align}
The corresponding compact matrix expressions are
\begin{align}
 \mathcal I_n^{ab;c}
 &=\big[\Gamma^{ca}K_2^b+K_2^a\Gamma^{cb}
        -2[V_F^c,K_1^a]K_2^b\big]_{nn},\label{eq:qgd-I-matrix}\\
 \mathcal J_n^{ab;c}
 &=\big[[K_1^c,V^a]K_2^b+K_2^a[K_1^c,V^b]\big]_{nn}.
 \label{eq:qgd-II-matrix}
\end{align}
No real part is taken until selecting the desired observable:
\begin{equation}
 \partial_cG_n^{ab}=\Rea\partial_cQ_n^{ab},\qquad
 \partial_c\Omega_n^{ab}=-2\Ima\partial_cQ_n^{ab}.
 \label{eq:qgd-observables}
\end{equation}
Equations~\eqref{eq:qgd-I}--\eqref{eq:qgd-II-matrix} retain the
one-sum/two-sum distinction in \code{QGD.tex} without repeating the intermediate bra-ket expansions.

\subsection{A complete quotient identity for the second derivative}
The QGQ expansions in the source notes are incomplete.
For a compact check of any completed expansion, define the phase-invariant pair
$Z_{nm}^{ab}=V_{nm}^aV_{mn}^b$ for $m\ne n$.
Differentiating $Q_n^{ab}=\sum_{m\ne n}Z_{nm}^{ab}/\Delta_{nm}^2$ twice gives
\begin{align}
 \partial_d\partial_c Q_n^{ab}=\sum_{m\ne n}\bigg[
 &\frac{\partial_d\partial_c Z_{nm}^{ab}}{\Delta_{nm}^2}
 -\frac{2\big[(\partial_c Z_{nm}^{ab})\Delta_{nm,d}
              +(\partial_d Z_{nm}^{ab})\Delta_{nm,c}
              +Z_{nm}^{ab}\Delta_{nm,cd}\big]}{\Delta_{nm}^3}
 \nonumber\\
 &+\frac{6Z_{nm}^{ab}\Delta_{nm,c}\Delta_{nm,d}}{\Delta_{nm}^4}\bigg].
 \label{eq:qgq-quotient}
\end{align}
In particular, the last coefficient is $+6$ in this scalar band-sum convention.
The numerator derivatives include changes of both eigenstates and operators.
For an arbitrary operator matrix $X$, they can be generated using
\begin{equation}
 X_{nm}^{;c}=\langle n|\partial_c\hat X|m\rangle+[K_1^c,X]_{nm},\qquad
 (K_q^a)_{nm}^{;c}=\frac{(V^a)_{nm}^{;c}}{\Delta_{nm}^q}
 -q\frac{V_{nm}^a\Delta_{nm,c}}{\Delta_{nm}^{q+1}},\quad n\ne m.
 \label{eq:qgd-recursion}
\end{equation}
Here $;c$ removes only eigenline phase connections, with the same choice as in
Sec.~\ref{sec:bcp}. Differentiating $\Gamma$ includes
$\langle n|\partial_c\partial_a\partial_b\mathcal H|m\rangle$.
It cannot be omitted for a general tight-binding Hamiltonian.
This identity is a complete differentiation rule, not a claim that the
partial QGQ matrix expansions in the source notes are complete.

\section{Energy derivatives and Drude building blocks}
\label{sec:energy}
Define the rotated bare derivatives
\begin{equation}
 V^a=\mathcal H^{,a},\quad W^{ab}=\mathcal H^{,ab},\quad
 Y^{abc}=\mathcal H^{,abc},\quad Z^{abcd}=\mathcal H^{,abcd}.
 \label{eq:energy-ingredients}
\end{equation}
Here $Z^{abcd}$ is a four-derivative Hamiltonian matrix, distinct from the
band-pair numerator $Z_{nm}^{ab}$ in Sec.~\ref{sec:qgd}.
For example,
\begin{align}
 V_{AA}^{a:b}&=W_{AA}^{ab}+V_{AB}^aD_{BA}^b-D_{AB}^bV_{BA}^a,
 \label{eq:invmass}\\
 W_{AA}^{ab:c}&=Y_{AA}^{abc}+W_{AB}^{ab}D_{BA}^c-D_{AB}^cW_{BA}^{ab},
 \label{eq:derW}\\
 V_{AA}^{a:bc}&=W_{AA}^{ab:c}
 +V_{AB}^{a:c}D_{BA}^b+V_{AB}^aD_{BA}^{b:c}
 -D_{AB}^{b:c}V_{BA}^a-D_{AB}^bV_{BA}^{a:c}.
 \label{eq:third-energy}
\end{align}
A further product derivative yields
\begin{align}
 V_{AA}^{a:bcd}={}&W_{AA}^{ab:cd}
 +V_{AB}^{a:cd}D_{BA}^b+V_{AB}^{a:c}D_{BA}^{b:d}
 +V_{AB}^{a:d}D_{BA}^{b:c}+V_{AB}^aD_{BA}^{b:cd}
 \nonumber\\
 &-D_{AB}^{b:cd}V_{BA}^a-D_{AB}^{b:c}V_{BA}^{a:d}
 -D_{AB}^{b:d}V_{BA}^{a:c}-D_{AB}^bV_{BA}^{a:cd}.
 \label{eq:fourth-energy}
\end{align}
For a smooth isolated group,
\begin{equation}
 \partial_a\partial_b\sum_{n\in A}\varepsilon_n=\Tr_A V^{a:b},\quad
 \partial_a\partial_b\partial_c\sum_{n\in A}\varepsilon_n=\Tr_A V^{a:bc},\quad
 \partial_a\partial_b\partial_c\partial_d\sum_{n\in A}\varepsilon_n=\Tr_A V^{a:bcd}.
 \label{eq:energy-traces}
\end{equation}
For a single isolated band the inverse effective mass is
$(m_n^{-1})^{ab}=\hbar^{-2}\partial_a\partial_b\varepsilon_n$.
The operator Hessian $W_{nn}^{ab}$ differs from the band-energy Hessian:
\begin{equation}
 \partial_a\partial_b\varepsilon_n
 =W_{nn}^{ab}+2\Rea\sum_{m\ne n}\frac{V_{nm}^aV_{mn}^b}{\Delta_{nm}}.
 \label{eq:hessian-distinction}
\end{equation}
Thus the orbital-quadrupole derivative of an operator Hessian must not be
replaced by \code{Der3E}. They correspond to different objects.

\section{Orbital magnetic quantities}
\label{sec:orbital}

\subsection{Orbital moment and zero-temperature magnetization}
With the charge convention in Eq.~\eqref{eq:conventions},
\begin{align}
 \widetilde m^c&=\frac{\ii e}{2\hbar}\epsilon_{cab}
                      (\widetilde G^{ab}-\widetilde H^{ab}),\label{eq:moment}\\
 \widetilde{\mathcal M}^c(\mu)
 &=-\frac{\ii e}{2\hbar}\epsilon_{cab}
    (\widetilde H^{ab}+\widetilde G^{ab}-2\mu\widetilde F^{ab}).
 \label{eq:magnetization}
\end{align}
The first is the symmetrized orbital-moment matrix; for an exactly degenerate
active group it reduces to the usual moment matrix.
For a single band,
$\mathcal M_n^c=m_n^c+(e/\hbar)(\mu-\varepsilon_n)\Omega_n^c$.
At zero temperature, using the occupied projector at each wave vector,
\begin{equation}
 M_{\mathrm{orb}}^c=\int[\dd\bm k]\Tr_{\mathrm{occ}}
                    \widetilde{\mathcal M}^c(\mu).
 \label{eq:magnetization-integral}
\end{equation}
At fixed $\mu$, the derivative of Eq.~\eqref{eq:magnetization} is
\begin{equation}
 (\widetilde{\mathcal M}^c)^{:d}
 =-\frac{\ii e}{2\hbar}\epsilon_{cab}
 (\widetilde H^{ab:d}+\widetilde G^{ab:d}-2\mu\widetilde F^{ab:d}).
 \label{eq:magnetization-derivative}
\end{equation}
Every projector derivative in $\widetilde H$ and $\widetilde G$ is included
through the interpolation and product rules. Keeping only
$\langle\partial_a u_m|Q(\partial_d\hat{\mathcal H})Q|\partial_b u_n\rangle$
is insufficient.

\subsection{Orbital magnetic quadrupole: four-term reference}
\label{sec:mqm}
To avoid the duplicated prefactors in \code{Wannier.tex}, this subsection
uses the $e=\hbar=1$ convention of Gao and Xiao~\cite{GaoXiao2018}.
It supplies the reduced kernel for comparison with the archive, not a new SI conversion.
For an isolated band define
\begin{align}
 \mathfrak m_{mn}^{j}
 &=\frac12\epsilon_{jab}\sum_{p\ne n}
       (v_{mp}^a+v_n^a\delta_{mp})A_{pn}^b,\label{eq:interband-moment}\\
 \Gamma_n^{i b}&=\langle u_n|\partial_i\partial_b\hat{\mathcal H}|u_n\rangle,
 \qquad g_n^{ib}\equiv G_n^{ib}.
\end{align}
Here $v^a=\partial_a\hat{\mathcal H}$ in this subsection's units.
Define four band kernels,
\begin{align}
 C_n^{ij}&=\frac23\Rea\sum_{m\ne n}A_{nm}^i\mathfrak m_{mn}^j,
 &I_n^{ij}&=-\frac1{12}\epsilon_{abj}\partial_a\Gamma_n^{ib},
 \nonumber\\
 L_n^{ij}&=-2\Rea\sum_{m\ne n}
                 \frac{A_{nm}^i\mathfrak m_{mn}^j}{\Delta_{nm}},
 &K_n^{ij}&=-\frac13\epsilon_{abj}\partial_aG_n^{ib}.
 \label{eq:mqm-kernels}
\end{align}
The quadrupole density is
\begin{equation}
 \mathcal Q^{ij}=\int[\dd\bm k]\sum_n
 \left\{f_n(C_n^{ij}+I_n^{ij})+\Phi_n(L_n^{ij}+K_n^{ij})\right\},
 \label{eq:mqm}
\end{equation}
where
\begin{equation}
 f_n=\frac1{1+e^{(\varepsilon_n-\mu)/(k_BT)}},\qquad
 \Phi_n=-k_BT\ln\left(1+e^{(\mu-\varepsilon_n)/(k_BT)}\right).
 \label{eq:grand-potential}
\end{equation}
At $T=0$, $f_n=\Theta(\mu-\varepsilon_n)$ and
$\Phi_n=(\varepsilon_n-\mu)f_n$.
The source-note labels ME1, W, ME2, QM correspond respectively to the
archived v1.8 labels CL, IMD, ME, QMD.

The physical operator Hessian in Eq.~\eqref{eq:mqm-kernels} equals the fixed-basis
matrix $W$ of Sec.~\ref{sec:energy} in that representation.
With a moving Wannier basis, operator derivatives and derivatives of
$\mathcal H^W$ require additional basis-connection terms.
This distinction must be checked when comparing the archived \code{DerWln}
implementation to the physical quadrupole formula.

The archived v1.8 calculator uses $(E-\mu)f$ for ME and QMD and defaults to
\code{constant_factor=1}. Its weighting agrees with
Eq.~\eqref{eq:grand-potential} at zero temperature, and does not by itself
implement the finite-temperature grand potential.
No additional $e/(2\hbar)$ factor should be inserted into
Eq.~\eqref{eq:mqm} while retaining the convention in Eq.~\eqref{eq:interband-moment}.
Conversion back to physical units must be derived for the full expression.

\section{Bridge to the v1.7 and v1.8 calculator catalogue}
\label{sec:catalogue}

\subsection{Fermi-sea and Fermi-surface kernels}
For a smooth periodic band quantity $X_n$ and vanishing boundary terms,
\begin{align}
 \int[\dd\bm k]\sum_n f_n\partial_aX_n
 &=-\int[\dd\bm k]\sum_n f'_n E_n^a X_n,\label{eq:sea-surface}\\
 \int[\dd\bm k]\sum_n f_n\partial_a\partial_dX_n
 &=-\int[\dd\bm k]\sum_n f'_n E_n^a\partial_dX_n,
 \label{eq:sea-surface2}
\end{align}
where $f'_n=\partial f/\partial\varepsilon|_{\varepsilon_n}$ is negative.
A surface kernel defined with $-f'$ has the same sign as the sea integral.
The code's occupation-derivative convention and outer factor must therefore
be checked together. At band crossings, use smooth group traces and the
matrix occupation rule rather than applying isolated-band formulae blindly.

\subsection{Tensor axes and formula-class index}
The following table records geometric kernel axes, excluding any leading
wave-vector and chemical-potential axes. It is a navigation aid for the
subsequent full version reconciliation.
\begin{center}
\setlength{\tabcolsep}{6pt}
\begin{tabular}{lll}
\toprule
Formula class & Geometric kernel & Raw Cartesian axes\\
\midrule
\code{Metric} & $G^{ab}$ or subspace metric & $(a,b)$\\
\code{DerMetric} & $G^{ab:d}$ & $(a,b,d)$\\
\code{BCP} & $\mathcal G^{ab}$ & $(a,b)$\\
\code{DerBCP} & $\mathcal G^{ab:d}$ & $(a,b,d)$\\
\code{Der2BCP} & $\mathcal G^{ab:de}$ & $(a,b,d,e)$\\
\code{DerOmega}, \code{Der2Omega} & $\Omega^{c:d}$, $\Omega^{c:de}$ & $(c,d)$, $(c,d,e)$\\
\code{VelBCP} & $E^c\mathcal G^{ab}$ & $(c,a,b)$\\
\code{VelDerBCP} & $E^c\mathcal G^{ab:d}$ & $(c,a,b,d)$\\
\code{VelDerOmega} & $E^a\Omega^{c:d}$ & $(a,c,d)$\\
\code{InvMass}, \code{Der3E}, \code{Der4E} & band-energy derivatives & derivative order\\
\code{Qorb_GaoXiao_*} & Eq.~\eqref{eq:mqm} & $(i,j)$\\
\bottomrule
\end{tabular}
\end{center}
For a common dipole convention $(d,a,b)$, transpose the three axes of
\code{DerBCP} from $(a,b,d)$ to $(d,a,b)$.
For a common quadrupole convention $(d,a,b,e)$, the corresponding permutation
of \code{Der2BCP} is $(2,0,1,3)$ in zero-based tensor-axis notation.
Only the trailing tensor axes should be permuted when leading axes are present.

\subsection{Items for the subsequent complete version reconciliation}
The archived source and \code{CHANGE_REPORT.md} identify the following scope:
\begin{itemize}
 \item v1.7 adds metric, BCP dipole/quadrupole, Berry-curvature quadrupole,
 fourth-energy-derivative, and early orbital-quadrupole calculators.
 \item v1.8 retains several formula classes but does not expose the complete
 v1.7 calculator set. It adds the four-term Gao--Xiao quadrupole and
 \code{Beta_GaoXiao} thermoelectric response.
 \item \code{Der2BCP_test} in v1.7 raises \code{NotImplementedError}.
 The production \code{Der2BCP} also carries a bug-related TODO in the source;
 a class name or a completed symbolic identity is not evidence of numerical validation.
 \item BCP denominator derivatives, active/complement choices, tensor axes,
 occupation signs, and SI prefactors need class-by-class comparison.
 Group-metric traces and sums of isolated-band metrics must be distinguished.
\end{itemize}
The thermoelectric response and the full nonlinear conductivity prefactors
are outside the completed source-note consolidation here.
They should be added as separate catalogue entries with their field convention,
thermal kernel, component symmetrization, and version-specific code mapping.

\appendix
\section{Corrections and retained evidence boundaries}
\label{app:corrections}
\begin{enumerate}
 \item \textbf{Connection and derivative notation.}
 The source notes used the same state notation for Bloch functions and
 columns of $U$. Equations~\eqref{eq:D}--\eqref{eq:M} distinguish them.
 The ordinary derivative $\partial_d X$ was also confused with the
 operator/comma derivative $X^{,d}$; Sec.~\ref{sec:derivative-conventions} repairs that distinction.
 The right-block gauge transformation has left index $n$, not the unrelated index $m$.
 \item \textbf{Curvature and Hermiticity.}
 The anti-Hermitian relation is $(D_{nm}^a)^*=-D_{mn}^a$.
 The $z$-curvature contraction uses $xy-yx$, not $xy-yz$.
 The internal trace curvature has the positive sign in
 Eq.~\eqref{eq:curvature-check}.
 The $+2\Ima Q$ convention in \code{QGD.tex} is replaced consistently by
 Eq.~\eqref{eq:metric-curvature}.
 \item \textbf{BCP factors and denominator derivatives.}
 The metric symmetrization and the symmetric energy-denominator construction
 each contribute a factor of $1/2$.
 An expanded derivative in \code{Wannier.tex} lost one of these factors
 and repeated terms instead of exchanging all $a,b$ indices.
 Equations~\eqref{eq:bcp-first}--\eqref{eq:bcp-sym} retain the two operations explicitly.
 The Sylvester identities supplement the single-band quotient rule without
 assuming diagonal intrablock velocity matrices.
 \item \textbf{Incomplete QGQ expansion.}
 \code{Wannier.tex} leaves its part III empty, and \code{QGD.tex} repeats
 the part-I label on a different group of terms.
 These partial expressions are not reproduced as a complete result.
 Equation~\eqref{eq:qgq-quotient} records the full quotient identity, including
 the $+6$ denominator term and the band-energy Hessian difference.
 \item \textbf{Orbital terms.}
 The energy-weighted Wannier matrix uses derivatives $a,b$, not $b,b$.
 Orbital-moment, magnetization, and interband magnetic-kernel symbols are separated.
 The original quadrupole expression mixed an outer $e/(2\hbar)$ with a magnetic
 matrix already containing that factor. Section~\ref{sec:mqm} uses a single
 explicit reduced-unit convention.
 \item \textbf{Exploratory derivations.}
 The unfinished ME1 projector manipulations, partial magnetization derivatives,
 and terminal Green-function identities are retained in the original notes.
 They are not needed for the reusable formula chain in this handbook.
 In particular, differentiation of a projected product must include all
 projector derivatives; differentiating only its explicit Hamiltonian is incomplete.
\end{enumerate}

\section{Consistency checks for later calculator validation}
\label{app:checks}
Useful algebraic checks include
\begin{align}
 Q_n^{ab}&=(Q_n^{ba})^*,\qquad G_n^{aa}\ge0,\nonumber\\
 \widetilde F_{mn}^{ab:s}&=(\widetilde F_{nm}^{ba:s})^*,\qquad
 \mathcal G_n^{ab}=\mathcal G_n^{ba},\nonumber\\
 \partial_d\partial_cX_n&=\partial_c\partial_dX_n
 \quad\text{for a smooth gauge-invariant scalar }X_n.
\end{align}
Compare the QGD band sums, matrix products, and centered finite differences
on a gapped multiband model. Test the BCP block denominator identities under
independent unitary rotations in both blocks.
For Brillouin-zone integrals, compare sea/surface kernels only after matching
axes, signs, units, occupations, and integration normalization.
These checks do not replace convergence tests on the eventual material model.

% Embedded bibliography keeps this working copy self-contained for the native
% single-file LaTeX compiler. Original notes/ref.bib is deliberately unchanged.
\begin{thebibliography}{9}
\bibitem{Liu2023CovariantderivativesBerry}
X. Liu, S. S. Tsirkin, and I. Souza,
Covariant derivatives of Berry-type quantities: Application to nonlinear transport,
\href{https://arxiv.org/abs/2303.10129}{arXiv:2303.10129} (2023).
\bibitem{GaoPRL2014FieldInducedPositional}
Y. Gao, S. A. Yang, and Q. Niu,
Field Induced Positional Shift of Bloch Electrons and Its Dynamical Implications,
\href{https://doi.org/10.1103/PhysRevLett.112.166601}{Phys. Rev. Lett. \textbf{112}, 166601} (2014).
\bibitem{GaoXiao2018}
Y. Gao and D. Xiao,
Orbital magnetic quadrupole moment and nonlinear anomalous thermoelectric transport,
\href{https://doi.org/10.1103/PhysRevB.98.060402}{Phys. Rev. B \textbf{98}, 060402(R)} (2018).
\end{thebibliography}
\end{document}
